\appendix
\section{Proofs}

\subsection{Proof of Lemma: QP Dual Problem}

\textit{Lemma.} The dual of \eqref{eq:qp_primal} is \eqref{eq:qp_dual}.

\textit{Proof.} The Lagrangian of \eqref{eq:qp_primal} is:
\begin{equation}
    \cL(d, \lambda) = \frac{1}{2} \|d + g_F\|^2 + \lambda^T (J_{verified} d)
\end{equation}
with $\lambda \geq 0$. Taking gradient with respect to $d$:
\begin{equation}
    \nabla_d \cL = d + g_F + J_{verified}^T \lambda = 0 \implies d^* = -g_F - J_{verified}^T \lambda
\end{equation}
Substituting into $\cL$:
\begin{align}
    \cL(\lambda) &= \frac{1}{2} \| -g_F - J_{verified}^T \lambda + g_F \|^2 + \lambda^T J_{verified}(-g_F - J_{verified}^T \lambda) \\
    &= \frac{1}{2} \lambda^T (J_{verified} J_{verified}^T) \lambda - \lambda^T J_{verified} g_F
\end{align}
which is exactly \eqref{eq:qp_dual}. \qed

\subsection{Proof of Property: Constraint Satisfaction}

\textit{Property.} The QP projection guarantees that for all verified simplices $j$:
\begin{equation}
    J_{verified}[j,:] \cdot d^* \leq 0
\end{equation}

\textit{Proof.} From the KKT conditions of the QP, complementary slackness requires:
\begin{equation}
    \lambda^*_j \cdot (J_{verified}[j,:] \cdot d^*) = 0 \quad \forall j
\end{equation}
Since $\lambda^* \geq 0$ and the primal constraint is $J_{verified} d^* \leq 0$, we have $J_{verified} d^* \leq 0$. \qed

\subsection{Proof of Property: Descent Property}

\textit{Property.} If $g_F \neq 0$ and the QP is feasible, the projected direction satisfies $g_F \cdot d^* \leq 0$.

\textit{Proof.} From $d^* = -g_F + v$ where $v = -J_{verified}^T \lambda^*$ (from the KKT conditions), we have:
\begin{equation}
    g_F \cdot d^* = -\|g_F\|^2 + g_F \cdot v \leq -\|g_F\|^2 + \|g_F\| \cdot \|v\|
\end{equation}
Since $v = -J_{verified}^T \lambda^*$ is a linear combination of rows of $J_{verified}^T$ with non-negative coefficients $\lambda^*$, and $J_{verified} d^* \leq 0$, we have $g_F \cdot v \leq 0$. Hence $g_F \cdot d^* \leq 0$. \qed

\section{Additional Experimental Details}

\subsection{Dataset and System Definitions}

The four dynamical systems are defined as follows:

\textbf{Simple2D:}
\begin{equation}
    \dot{x}_1 = x_2, \quad \dot{x}_2 = -x_1 - x_2 + \frac{1}{3}x_1^3
\end{equation}

\textbf{Barrier1:} Safety set $h(x) = 1 - x_1^2 - x_2^2 \geq 0$

\textbf{Barrier2:} Safety set defined by multiple polynomial constraints

\textbf{Barrier3:} Safety set with complex geometry including non-convex regions

\subsection{Hyperparameters}

Table~\ref{tab:hyperparams} lists the key hyperparameters used in our experiments.

\begin{table}[t]
    \caption{Hyperparameters}
    \label{tab:hyperparams}
    \centering
    \begin{tabular}{lc}
        \toprule
        Parameter & Value \\
        \midrule
        RS Jacobian samples ($N$) & 100 \\
        RS smoothing ($\sigma$) & 0.01 \\
        Inner loop steps & 5 \\
        Learning rate (depth 10) & $5 \times 10^{-3}$ \\
        Learning rate (depth 12) & $1 \times 10^{-3}$ \\
        Learning rate (depth 15) & $5 \times 10^{-4}$ \\
        Depth schedule & $\{10, 12, 15, 18, 20\}$ \\
        Plateau threshold ($\tau$) & 0.5\% \\
        Max stagnant iterations & 5 \\
        \bottomrule
    \end{tabular}
\end{table}